\section{Discussion}
\label{sec:discussion}

\subsection{Main empirical takeaways}

The Qwen2.5 and Qwen3 evidence supports three robust observations.
First, on RTX4090, speculative decoding can provide substantial acceleration,
but not for every setting; high-$k$ stochastic can drop below baseline.
Second, speedup declines as $k$ increases, even when accepted block size grows.
Third, deterministic decoding consistently outperforms stochastic decoding
for matched $k$ in both run families.

The cross-hardware extension adds an important qualification.
While policy ordering trends transfer from RTX4090 to RTX5090-laptop,
absolute speedup magnitudes do not.
This supports a practical interpretation: generalization claims should emphasize
portable trends and decision rules, not fixed throughput numbers.
The adaptive-$k$ run set was evaluated at a single fixed rescue-threshold,
window, and cooling configuration in this study, reinforcing its framing as a
stability-oriented mechanism rather than a tuned throughput optimizer; a
systematic sweep over these hyperparameters is identified as future work in
Section~\ref{sec:threats}.

\subsection{Novelty and contribution scope}

This paper's contribution is primarily empirical and systems-oriented.
We do not claim a new speculative decoding algorithm that universally
improves throughput.
Instead, we provide a controlled, reproducible characterization that links
acceptance dynamics, token timing, and latency tails to practical operating
choices on consumer hardware.

\subsection{Variance and evaluation depth}

To strengthen rigor, we report sample-level mean$\pm$std summaries,
time to first token (TTFT), time per output token (TPOT), and
Holm-Bonferroni-corrected paired $t$-tests with Cohen's $d$ effect sizes.
These additions improve variability reporting and statistical clarity, and
show that five of six fixed-$k$ configurations achieve significant,
large-effect speedup, while stochastic $k{=}16$ does not.
We further add seed-stability extensions for both Qwen2.5 and Qwen3 on
RTX4090 and RTX5090-laptop at $k{=}4$ (seeds $123$ and $999$ in addition to
the primary seed $42$ where available). They preserve
deterministic-over-stochastic ordering and show small seed-to-seed drift in
acceptance and token timing at this operating point on both devices.
We also expand quality reporting beyond a single summarization metric
(ROUGE-L, BLEU, and BERTScore, computed on the real CNN/DailyMail outputs).
All three metrics move together and show a genuine quality cost under
speculative decoding relative to the autoregressive baseline, largest at
$k{=}4$ and partially recovering at higher $k$ -- the reverse of the speed
trend, indicating a real speed--quality trade-off rather than a free
lunch at low $k$. We acknowledge that none of these metrics fully capture
factuality or task-specific utility.

\section{Threats to validity}
\label{sec:threats}

\textbf{Hardware and runtime dependence.}
Results remain sensitive to implementation stack and runtime path.
The RTX4090 and RTX5090-laptop comparison demonstrates that ranking trends can
transfer while absolute gains diverge.

\textbf{Model-family scope.}
This study centers on same-family Qwen2.5 and Qwen3 target/draft pairings.
Cross-family speculative pairings may show different acceptance and quality
behavior.

\textbf{Evaluation breadth.}
The benchmark suite covers reasoning, knowledge QA, and summarization, but does
not exhaust long-context, safety, or factual consistency requirements.

\textbf{Sample-level versus seed-level dispersion.}
The mean$\pm$std figures reported in Section~\ref{sec:results} reflect
per-sample dispersion within each evaluation seed. We now include repeated
seed runs (123 and 999) for the Qwen2.5 and Qwen3 $k{=}4$ deterministic and
stochastic settings on RTX4090 and RTX5090-laptop, providing initial
seed-level stability evidence across both model families and devices.
However, full multi-seed repetition across the entire
$k\in\{4,8,16\}$ grid was not performed and remains future work.

\textbf{Adaptive-$k$ hyperparameter scope.}
The adaptive-$k$ run set was evaluated at a single fixed rescue-threshold,
window-size, and cooling-period configuration. A systematic ablation across
these hyperparameters was not performed and is left for future work.

\section{Conclusion}
\label{sec:conclusion}

This paper presents reproducible deployment guidance for
consumer GPUs using same-family speculative decoding in both Qwen2.5 and
Qwen3 settings. On RTX4090, deterministic $k{=}4$ remains the strongest
operating point family-wise, with monotonic degradation as $k$ increases.
Cross-hardware validation on RTX5090-laptop preserves the same policy ordering
trend but not the same absolute speedups, emphasizing that portability must be
interpreted at the trend level.

The adaptive-$k$ run set remains useful as a runtime-stability policy layer,
but is not presented as a universal throughput optimizer.
Future work will expand hardware diversity, include additional model families,
and evaluate longer-context and richer quality diagnostics.
